% Options for packages loaded elsewhere
\PassOptionsToPackage{unicode}{hyperref}
\PassOptionsToPackage{hyphens}{url}
%
\documentclass[
]{article}
\usepackage{amsmath,amssymb}
\usepackage{iftex}
\ifPDFTeX
  \usepackage[T1]{fontenc}
  \usepackage[utf8]{inputenc}
  \usepackage{textcomp} % provide euro and other symbols
\else % if luatex or xetex
  \usepackage{unicode-math} % this also loads fontspec
  \defaultfontfeatures{Scale=MatchLowercase}
  \defaultfontfeatures[\rmfamily]{Ligatures=TeX,Scale=1}
\fi
\usepackage{lmodern}
\ifPDFTeX\else
  % xetex/luatex font selection
\fi
% Use upquote if available, for straight quotes in verbatim environments
\IfFileExists{upquote.sty}{\usepackage{upquote}}{}
\IfFileExists{microtype.sty}{% use microtype if available
  \usepackage[]{microtype}
  \UseMicrotypeSet[protrusion]{basicmath} % disable protrusion for tt fonts
}{}
\makeatletter
\@ifundefined{KOMAClassName}{% if non-KOMA class
  \IfFileExists{parskip.sty}{%
    \usepackage{parskip}
  }{% else
    \setlength{\parindent}{0pt}
    \setlength{\parskip}{6pt plus 2pt minus 1pt}}
}{% if KOMA class
  \KOMAoptions{parskip=half}}
\makeatother
\usepackage{xcolor}
\usepackage[margin=1in]{geometry}
\usepackage{color}
\usepackage{fancyvrb}
\newcommand{\VerbBar}{|}
\newcommand{\VERB}{\Verb[commandchars=\\\{\}]}
\DefineVerbatimEnvironment{Highlighting}{Verbatim}{commandchars=\\\{\}}
% Add ',fontsize=\small' for more characters per line
\usepackage{framed}
\definecolor{shadecolor}{RGB}{248,248,248}
\newenvironment{Shaded}{\begin{snugshade}}{\end{snugshade}}
\newcommand{\AlertTok}[1]{\textcolor[rgb]{0.94,0.16,0.16}{#1}}
\newcommand{\AnnotationTok}[1]{\textcolor[rgb]{0.56,0.35,0.01}{\textbf{\textit{#1}}}}
\newcommand{\AttributeTok}[1]{\textcolor[rgb]{0.13,0.29,0.53}{#1}}
\newcommand{\BaseNTok}[1]{\textcolor[rgb]{0.00,0.00,0.81}{#1}}
\newcommand{\BuiltInTok}[1]{#1}
\newcommand{\CharTok}[1]{\textcolor[rgb]{0.31,0.60,0.02}{#1}}
\newcommand{\CommentTok}[1]{\textcolor[rgb]{0.56,0.35,0.01}{\textit{#1}}}
\newcommand{\CommentVarTok}[1]{\textcolor[rgb]{0.56,0.35,0.01}{\textbf{\textit{#1}}}}
\newcommand{\ConstantTok}[1]{\textcolor[rgb]{0.56,0.35,0.01}{#1}}
\newcommand{\ControlFlowTok}[1]{\textcolor[rgb]{0.13,0.29,0.53}{\textbf{#1}}}
\newcommand{\DataTypeTok}[1]{\textcolor[rgb]{0.13,0.29,0.53}{#1}}
\newcommand{\DecValTok}[1]{\textcolor[rgb]{0.00,0.00,0.81}{#1}}
\newcommand{\DocumentationTok}[1]{\textcolor[rgb]{0.56,0.35,0.01}{\textbf{\textit{#1}}}}
\newcommand{\ErrorTok}[1]{\textcolor[rgb]{0.64,0.00,0.00}{\textbf{#1}}}
\newcommand{\ExtensionTok}[1]{#1}
\newcommand{\FloatTok}[1]{\textcolor[rgb]{0.00,0.00,0.81}{#1}}
\newcommand{\FunctionTok}[1]{\textcolor[rgb]{0.13,0.29,0.53}{\textbf{#1}}}
\newcommand{\ImportTok}[1]{#1}
\newcommand{\InformationTok}[1]{\textcolor[rgb]{0.56,0.35,0.01}{\textbf{\textit{#1}}}}
\newcommand{\KeywordTok}[1]{\textcolor[rgb]{0.13,0.29,0.53}{\textbf{#1}}}
\newcommand{\NormalTok}[1]{#1}
\newcommand{\OperatorTok}[1]{\textcolor[rgb]{0.81,0.36,0.00}{\textbf{#1}}}
\newcommand{\OtherTok}[1]{\textcolor[rgb]{0.56,0.35,0.01}{#1}}
\newcommand{\PreprocessorTok}[1]{\textcolor[rgb]{0.56,0.35,0.01}{\textit{#1}}}
\newcommand{\RegionMarkerTok}[1]{#1}
\newcommand{\SpecialCharTok}[1]{\textcolor[rgb]{0.81,0.36,0.00}{\textbf{#1}}}
\newcommand{\SpecialStringTok}[1]{\textcolor[rgb]{0.31,0.60,0.02}{#1}}
\newcommand{\StringTok}[1]{\textcolor[rgb]{0.31,0.60,0.02}{#1}}
\newcommand{\VariableTok}[1]{\textcolor[rgb]{0.00,0.00,0.00}{#1}}
\newcommand{\VerbatimStringTok}[1]{\textcolor[rgb]{0.31,0.60,0.02}{#1}}
\newcommand{\WarningTok}[1]{\textcolor[rgb]{0.56,0.35,0.01}{\textbf{\textit{#1}}}}
\usepackage{graphicx}
\makeatletter
\newsavebox\pandoc@box
\newcommand*\pandocbounded[1]{% scales image to fit in text height/width
  \sbox\pandoc@box{#1}%
  \Gscale@div\@tempa{\textheight}{\dimexpr\ht\pandoc@box+\dp\pandoc@box\relax}%
  \Gscale@div\@tempb{\linewidth}{\wd\pandoc@box}%
  \ifdim\@tempb\p@<\@tempa\p@\let\@tempa\@tempb\fi% select the smaller of both
  \ifdim\@tempa\p@<\p@\scalebox{\@tempa}{\usebox\pandoc@box}%
  \else\usebox{\pandoc@box}%
  \fi%
}
% Set default figure placement to htbp
\def\fps@figure{htbp}
\makeatother
\setlength{\emergencystretch}{3em} % prevent overfull lines
\providecommand{\tightlist}{%
  \setlength{\itemsep}{0pt}\setlength{\parskip}{0pt}}
\setcounter{secnumdepth}{-\maxdimen} % remove section numbering
\usepackage{graphicx}
\usepackage{float}
\usepackage{tikz}
\usepackage{xcolor}
\usepackage{bookmark}
\IfFileExists{xurl.sty}{\usepackage{xurl}}{} % add URL line breaks if available
\urlstyle{same}
\hypersetup{
  hidelinks,
  pdfcreator={LaTeX via pandoc}}

\author{}
\date{\vspace{-2.5em}}

\begin{document}

\begin{Shaded}
\begin{Highlighting}[]
\FunctionTok{options}\NormalTok{(}\AttributeTok{OutDec =} \StringTok{"."}\NormalTok{)  }\CommentTok{\# Asegurar punto decimal en este chunk}

\CommentTok{\# Establecer semilla aleatoria}
\FunctionTok{set.seed}\NormalTok{(}\FunctionTok{sample}\NormalTok{(}\DecValTok{1}\SpecialCharTok{:}\DecValTok{10000}\NormalTok{, }\DecValTok{1}\NormalTok{))}

\CommentTok{\# Aleatorización del contexto del problema}
\NormalTok{contextos }\OtherTok{\textless{}{-}} \FunctionTok{c}\NormalTok{(}
  \StringTok{"ciudad"}\NormalTok{, }\StringTok{"municipio"}\NormalTok{, }\StringTok{"localidad"}\NormalTok{, }\StringTok{"comunidad"}\NormalTok{, }\StringTok{"pueblo"}\NormalTok{,}
  \StringTok{"distrito"}\NormalTok{, }\StringTok{"barrio"}\NormalTok{, }\StringTok{"región"}\NormalTok{, }\StringTok{"provincia"}\NormalTok{, }\StringTok{"zona"}
\NormalTok{)}
\NormalTok{contexto }\OtherTok{\textless{}{-}} \FunctionTok{sample}\NormalTok{(contextos, }\DecValTok{1}\NormalTok{)}

\CommentTok{\# Aleatorización del tipo de competencia}
\NormalTok{competencias }\OtherTok{\textless{}{-}} \FunctionTok{c}\NormalTok{(}
  \StringTok{"maratón"}\NormalTok{, }\StringTok{"carrera"}\NormalTok{, }\StringTok{"competencia atlética"}\NormalTok{, }\StringTok{"torneo deportivo"}\NormalTok{,}
  \StringTok{"olimpiada deportiva"}\NormalTok{, }\StringTok{"evento deportivo"}\NormalTok{, }\StringTok{"justa deportiva"}\NormalTok{,}
  \StringTok{"campeonato"}\NormalTok{, }\StringTok{"concurso deportivo"}\NormalTok{, }\StringTok{"prueba atlética"}
\NormalTok{)}
\NormalTok{competencia }\OtherTok{\textless{}{-}} \FunctionTok{sample}\NormalTok{(competencias, }\DecValTok{1}\NormalTok{)}

\CommentTok{\# Aleatorización del grupo de edad}
\NormalTok{edades }\OtherTok{\textless{}{-}} \FunctionTok{c}\NormalTok{(}
  \StringTok{"menores de 15 años"}\NormalTok{, }\StringTok{"menores de 16 años"}\NormalTok{, }\StringTok{"menores de 14 años"}\NormalTok{,}
  \StringTok{"niños y niñas de 10 a 15 años"}\NormalTok{, }\StringTok{"jóvenes de 12 a 15 años"}\NormalTok{,}
  \StringTok{"estudiantes de primaria y secundaria"}\NormalTok{, }\StringTok{"categoría infantil"}\NormalTok{,}
  \StringTok{"categoría juvenil"}\NormalTok{, }\StringTok{"niños y adolescentes"}\NormalTok{, }\StringTok{"estudiantes menores de edad"}
\NormalTok{)}
\NormalTok{grupo\_edad }\OtherTok{\textless{}{-}} \FunctionTok{sample}\NormalTok{(edades, }\DecValTok{1}\NormalTok{)}

\CommentTok{\# Aleatorización del premio total (en millones)}
\CommentTok{\# Generamos valores que sean divisibles por 10 para facilitar cálculos}
\NormalTok{premios\_posibles }\OtherTok{\textless{}{-}} \FunctionTok{c}\NormalTok{(}\DecValTok{30}\NormalTok{, }\DecValTok{40}\NormalTok{, }\DecValTok{50}\NormalTok{, }\DecValTok{60}\NormalTok{, }\DecValTok{70}\NormalTok{, }\DecValTok{80}\NormalTok{, }\DecValTok{90}\NormalTok{, }\DecValTok{100}\NormalTok{, }\DecValTok{120}\NormalTok{, }\DecValTok{150}\NormalTok{)}
\NormalTok{premio\_total }\OtherTok{\textless{}{-}} \FunctionTok{sample}\NormalTok{(premios\_posibles, }\DecValTok{1}\NormalTok{)}

\CommentTok{\# Aleatorización de términos para el enunciado}
\NormalTok{terminos\_premiar }\OtherTok{\textless{}{-}} \FunctionTok{c}\NormalTok{(}\StringTok{"premiar"}\NormalTok{, }\StringTok{"recompensar"}\NormalTok{, }\StringTok{"reconocer"}\NormalTok{, }\StringTok{"galardonar"}\NormalTok{, }\StringTok{"incentivar"}\NormalTok{)}
\NormalTok{termino\_premiar }\OtherTok{\textless{}{-}} \FunctionTok{sample}\NormalTok{(terminos\_premiar, }\DecValTok{1}\NormalTok{)}

\NormalTok{terminos\_participantes }\OtherTok{\textless{}{-}} \FunctionTok{c}\NormalTok{(}\StringTok{"participantes"}\NormalTok{, }\StringTok{"competidores"}\NormalTok{, }\StringTok{"concursantes"}\NormalTok{, }\StringTok{"atletas"}\NormalTok{, }\StringTok{"deportistas"}\NormalTok{)}
\NormalTok{termino\_participantes }\OtherTok{\textless{}{-}} \FunctionTok{sample}\NormalTok{(terminos\_participantes, }\DecValTok{1}\NormalTok{)}

\NormalTok{terminos\_cuenta }\OtherTok{\textless{}{-}} \FunctionTok{c}\NormalTok{(}\StringTok{"cuenta con"}\NormalTok{, }\StringTok{"dispone de"}\NormalTok{, }\StringTok{"tiene asignado"}\NormalTok{, }\StringTok{"ha destinado"}\NormalTok{, }\StringTok{"ha reservado"}\NormalTok{)}
\NormalTok{termino\_cuenta }\OtherTok{\textless{}{-}} \FunctionTok{sample}\NormalTok{(terminos\_cuenta, }\DecValTok{1}\NormalTok{)}

\NormalTok{terminos\_repartiran }\OtherTok{\textless{}{-}} \FunctionTok{c}\NormalTok{(}\StringTok{"repartirán"}\NormalTok{, }\StringTok{"distribuirán"}\NormalTok{, }\StringTok{"dividirán"}\NormalTok{, }\StringTok{"asignarán"}\NormalTok{, }\StringTok{"otorgarán"}\NormalTok{)}
\NormalTok{termino\_repartiran }\OtherTok{\textless{}{-}} \FunctionTok{sample}\NormalTok{(terminos\_repartiran, }\DecValTok{1}\NormalTok{)}

\NormalTok{terminos\_puestos }\OtherTok{\textless{}{-}} \FunctionTok{c}\NormalTok{(}\StringTok{"primeros puestos"}\NormalTok{, }\StringTok{"primeras posiciones"}\NormalTok{, }\StringTok{"ganadores"}\NormalTok{, }\StringTok{"mejores lugares"}\NormalTok{, }\StringTok{"primeros lugares"}\NormalTok{)}
\NormalTok{termino\_puestos }\OtherTok{\textless{}{-}} \FunctionTok{sample}\NormalTok{(terminos\_puestos, }\DecValTok{1}\NormalTok{)}

\NormalTok{terminos\_dinero }\OtherTok{\textless{}{-}} \FunctionTok{c}\NormalTok{(}\StringTok{"dinero"}\NormalTok{, }\StringTok{"premio"}\NormalTok{, }\StringTok{"recompensa"}\NormalTok{, }\StringTok{"incentivo"}\NormalTok{, }\StringTok{"monto"}\NormalTok{)}
\NormalTok{termino\_dinero }\OtherTok{\textless{}{-}} \FunctionTok{sample}\NormalTok{(terminos\_dinero, }\DecValTok{1}\NormalTok{)}

\CommentTok{\# Aleatorización de fracciones para los puestos}
\CommentTok{\# Definimos conjuntos de fracciones que sumen menos de 1 para que quede algo para el tercer puesto}
\NormalTok{conjuntos\_fracciones }\OtherTok{\textless{}{-}} \FunctionTok{list}\NormalTok{(}
  \FunctionTok{c}\NormalTok{(}\StringTok{"1/2"}\NormalTok{, }\StringTok{"2/5"}\NormalTok{),  }\CommentTok{\# Suma 9/10, queda 1/10}
  \FunctionTok{c}\NormalTok{(}\StringTok{"1/3"}\NormalTok{, }\StringTok{"1/2"}\NormalTok{),  }\CommentTok{\# Suma 5/6, queda 1/6}
  \FunctionTok{c}\NormalTok{(}\StringTok{"2/5"}\NormalTok{, }\StringTok{"1/2"}\NormalTok{),  }\CommentTok{\# Suma 9/10, queda 1/10}
  \FunctionTok{c}\NormalTok{(}\StringTok{"3/5"}\NormalTok{, }\StringTok{"1/4"}\NormalTok{),  }\CommentTok{\# Suma 17/20, queda 3/20}
  \FunctionTok{c}\NormalTok{(}\StringTok{"1/2"}\NormalTok{, }\StringTok{"1/3"}\NormalTok{),  }\CommentTok{\# Suma 5/6, queda 1/6}
  \FunctionTok{c}\NormalTok{(}\StringTok{"3/4"}\NormalTok{, }\StringTok{"1/8"}\NormalTok{),  }\CommentTok{\# Suma 7/8, queda 1/8}
  \FunctionTok{c}\NormalTok{(}\StringTok{"2/3"}\NormalTok{, }\StringTok{"1/5"}\NormalTok{),  }\CommentTok{\# Suma 13/15, queda 2/15}
  \FunctionTok{c}\NormalTok{(}\StringTok{"3/5"}\NormalTok{, }\StringTok{"1/3"}\NormalTok{),  }\CommentTok{\# Suma 14/15, queda 1/15}
  \FunctionTok{c}\NormalTok{(}\StringTok{"1/2"}\NormalTok{, }\StringTok{"3/8"}\NormalTok{),  }\CommentTok{\# Suma 7/8, queda 1/8}
  \FunctionTok{c}\NormalTok{(}\StringTok{"3/5"}\NormalTok{, }\StringTok{"3/10"}\NormalTok{)  }\CommentTok{\# Suma 9/10, queda 1/10}
\NormalTok{)}

\CommentTok{\# Seleccionar un conjunto aleatorio de fracciones}
\NormalTok{indice\_conjunto }\OtherTok{\textless{}{-}} \FunctionTok{sample}\NormalTok{(}\DecValTok{1}\SpecialCharTok{:}\FunctionTok{length}\NormalTok{(conjuntos\_fracciones), }\DecValTok{1}\NormalTok{)}
\NormalTok{fracciones\_seleccionadas }\OtherTok{\textless{}{-}}\NormalTok{ conjuntos\_fracciones[[indice\_conjunto]]}

\CommentTok{\# Asignar fracciones a los puestos}
\NormalTok{fraccion\_primer\_puesto }\OtherTok{\textless{}{-}}\NormalTok{ fracciones\_seleccionadas[}\DecValTok{1}\NormalTok{]}
\NormalTok{fraccion\_segundo\_puesto }\OtherTok{\textless{}{-}}\NormalTok{ fracciones\_seleccionadas[}\DecValTok{2}\NormalTok{]}

\CommentTok{\# Convertir fracciones a valores numéricos para cálculos}
\NormalTok{convertir\_fraccion }\OtherTok{\textless{}{-}} \ControlFlowTok{function}\NormalTok{(fraccion) \{}
\NormalTok{  partes }\OtherTok{\textless{}{-}} \FunctionTok{strsplit}\NormalTok{(fraccion, }\StringTok{"/"}\NormalTok{)[[}\DecValTok{1}\NormalTok{]]}
  \FunctionTok{return}\NormalTok{(}\FunctionTok{as.numeric}\NormalTok{(partes[}\DecValTok{1}\NormalTok{]) }\SpecialCharTok{/} \FunctionTok{as.numeric}\NormalTok{(partes[}\DecValTok{2}\NormalTok{]))}
\NormalTok{\}}

\NormalTok{valor\_primer\_puesto }\OtherTok{\textless{}{-}} \FunctionTok{convertir\_fraccion}\NormalTok{(fraccion\_primer\_puesto)}
\NormalTok{valor\_segundo\_puesto }\OtherTok{\textless{}{-}} \FunctionTok{convertir\_fraccion}\NormalTok{(fraccion\_segundo\_puesto)}

\CommentTok{\# Calcular la fracción y el valor para el tercer puesto}
\NormalTok{valor\_tercer\_puesto }\OtherTok{\textless{}{-}} \DecValTok{1} \SpecialCharTok{{-}}\NormalTok{ (valor\_primer\_puesto }\SpecialCharTok{+}\NormalTok{ valor\_segundo\_puesto)}

\CommentTok{\# Verificar que el valor del tercer puesto sea positivo}
\FunctionTok{test\_that}\NormalTok{(}\StringTok{"El valor del tercer puesto es positivo"}\NormalTok{, \{}
  \FunctionTok{expect\_true}\NormalTok{(valor\_tercer\_puesto }\SpecialCharTok{\textgreater{}} \DecValTok{0}\NormalTok{)}
\NormalTok{\})}
\end{Highlighting}
\end{Shaded}

Test passed

\begin{Shaded}
\begin{Highlighting}[]
\CommentTok{\# Calcular los montos en millones para cada puesto}
\NormalTok{monto\_primer\_puesto }\OtherTok{\textless{}{-}}\NormalTok{ premio\_total }\SpecialCharTok{*}\NormalTok{ valor\_primer\_puesto}
\NormalTok{monto\_segundo\_puesto }\OtherTok{\textless{}{-}}\NormalTok{ premio\_total }\SpecialCharTok{*}\NormalTok{ valor\_segundo\_puesto}
\NormalTok{monto\_tercer\_puesto }\OtherTok{\textless{}{-}}\NormalTok{ premio\_total }\SpecialCharTok{*}\NormalTok{ valor\_tercer\_puesto}

\CommentTok{\# Redondear a números enteros si es necesario}
\NormalTok{monto\_primer\_puesto }\OtherTok{\textless{}{-}} \FunctionTok{round}\NormalTok{(monto\_primer\_puesto)}
\NormalTok{monto\_segundo\_puesto }\OtherTok{\textless{}{-}} \FunctionTok{round}\NormalTok{(monto\_segundo\_puesto)}
\NormalTok{monto\_tercer\_puesto }\OtherTok{\textless{}{-}} \FunctionTok{round}\NormalTok{(monto\_tercer\_puesto)}

\CommentTok{\# Ajustar el tercer puesto para asegurar que la suma sea exactamente el premio total}
\NormalTok{suma\_actual }\OtherTok{\textless{}{-}}\NormalTok{ monto\_primer\_puesto }\SpecialCharTok{+}\NormalTok{ monto\_segundo\_puesto }\SpecialCharTok{+}\NormalTok{ monto\_tercer\_puesto}
\ControlFlowTok{if}\NormalTok{ (suma\_actual }\SpecialCharTok{!=}\NormalTok{ premio\_total) \{}
\NormalTok{  monto\_tercer\_puesto }\OtherTok{\textless{}{-}}\NormalTok{ premio\_total }\SpecialCharTok{{-}}\NormalTok{ (monto\_primer\_puesto }\SpecialCharTok{+}\NormalTok{ monto\_segundo\_puesto)}
\NormalTok{\}}

\CommentTok{\# Verificar que la suma de los tres montos sea igual al premio total}
\FunctionTok{test\_that}\NormalTok{(}\StringTok{"La suma de los tres montos es igual al premio total"}\NormalTok{, \{}
  \FunctionTok{expect\_equal}\NormalTok{(monto\_primer\_puesto }\SpecialCharTok{+}\NormalTok{ monto\_segundo\_puesto }\SpecialCharTok{+}\NormalTok{ monto\_tercer\_puesto, premio\_total)}
\NormalTok{\})}
\end{Highlighting}
\end{Shaded}

Test passed

\begin{Shaded}
\begin{Highlighting}[]
\CommentTok{\# Generar opciones de respuesta}
\CommentTok{\# La respuesta correcta es el monto del tercer puesto}
\NormalTok{respuesta\_correcta }\OtherTok{\textless{}{-}}\NormalTok{ monto\_tercer\_puesto}

\CommentTok{\# Generar distractores plausibles}
\NormalTok{distractor1 }\OtherTok{\textless{}{-}} \FunctionTok{round}\NormalTok{(premio\_total }\SpecialCharTok{*} \FloatTok{0.05}\NormalTok{)  }\CommentTok{\# 5\% del premio total}
\NormalTok{distractor2 }\OtherTok{\textless{}{-}} \FunctionTok{round}\NormalTok{(premio\_total }\SpecialCharTok{*} \FloatTok{0.2}\NormalTok{)   }\CommentTok{\# 20\% del premio total}
\NormalTok{distractor3 }\OtherTok{\textless{}{-}} \FunctionTok{round}\NormalTok{(premio\_total }\SpecialCharTok{*} \FloatTok{0.3}\NormalTok{)   }\CommentTok{\# 30\% del premio total}

\CommentTok{\# Asegurarse de que todos los distractores son diferentes de la respuesta correcta}
\ControlFlowTok{if}\NormalTok{ (distractor1 }\SpecialCharTok{==}\NormalTok{ respuesta\_correcta) distractor1 }\OtherTok{\textless{}{-}}\NormalTok{ distractor1 }\SpecialCharTok{+} \DecValTok{1}
\ControlFlowTok{if}\NormalTok{ (distractor2 }\SpecialCharTok{==}\NormalTok{ respuesta\_correcta) distractor2 }\OtherTok{\textless{}{-}}\NormalTok{ distractor2 }\SpecialCharTok{+} \DecValTok{2}
\ControlFlowTok{if}\NormalTok{ (distractor3 }\SpecialCharTok{==}\NormalTok{ respuesta\_correcta) distractor3 }\OtherTok{\textless{}{-}}\NormalTok{ distractor3 }\SpecialCharTok{{-}} \DecValTok{2}

\CommentTok{\# Asegurarse de que todos los distractores son diferentes entre sí}
\ControlFlowTok{while}\NormalTok{ (}\FunctionTok{length}\NormalTok{(}\FunctionTok{unique}\NormalTok{(}\FunctionTok{c}\NormalTok{(distractor1, distractor2, distractor3))) }\SpecialCharTok{\textless{}} \DecValTok{3}\NormalTok{) \{}
  \ControlFlowTok{if}\NormalTok{ (distractor1 }\SpecialCharTok{==}\NormalTok{ distractor2) distractor1 }\OtherTok{\textless{}{-}}\NormalTok{ distractor1 }\SpecialCharTok{+} \DecValTok{1}
  \ControlFlowTok{if}\NormalTok{ (distractor2 }\SpecialCharTok{==}\NormalTok{ distractor3) distractor2 }\OtherTok{\textless{}{-}}\NormalTok{ distractor2 }\SpecialCharTok{+} \DecValTok{2}
  \ControlFlowTok{if}\NormalTok{ (distractor1 }\SpecialCharTok{==}\NormalTok{ distractor3) distractor3 }\OtherTok{\textless{}{-}}\NormalTok{ distractor3 }\SpecialCharTok{+} \DecValTok{3}
\NormalTok{\}}

\CommentTok{\# Crear un vector con todas las opciones y mezclarlas}
\NormalTok{opciones }\OtherTok{\textless{}{-}} \FunctionTok{c}\NormalTok{(respuesta\_correcta, distractor1, distractor2, distractor3)}
\FunctionTok{names}\NormalTok{(opciones) }\OtherTok{\textless{}{-}} \FunctionTok{c}\NormalTok{(}\StringTok{"correcta"}\NormalTok{, }\StringTok{"distractor1"}\NormalTok{, }\StringTok{"distractor2"}\NormalTok{, }\StringTok{"distractor3"}\NormalTok{)}
\NormalTok{opciones\_mezcladas }\OtherTok{\textless{}{-}} \FunctionTok{sample}\NormalTok{(opciones)}

\CommentTok{\# Identificar la posición de la respuesta correcta en las opciones mezcladas}
\NormalTok{indice\_correcto }\OtherTok{\textless{}{-}} \FunctionTok{which}\NormalTok{(opciones\_mezcladas }\SpecialCharTok{==}\NormalTok{ respuesta\_correcta)}

\CommentTok{\# Crear el vector de solución para r{-}exams}
\NormalTok{solucion }\OtherTok{\textless{}{-}} \FunctionTok{rep}\NormalTok{(}\DecValTok{0}\NormalTok{, }\DecValTok{4}\NormalTok{)}
\NormalTok{solucion[indice\_correcto] }\OtherTok{\textless{}{-}} \DecValTok{1}

\CommentTok{\# Aleatorización de colores para la tabla TikZ}
\NormalTok{paletas\_colores }\OtherTok{\textless{}{-}} \FunctionTok{list}\NormalTok{(}
  \FunctionTok{c}\NormalTok{(}\StringTok{"\#4285F4"}\NormalTok{, }\StringTok{"\#EA4335"}\NormalTok{, }\StringTok{"\#FBBC05"}\NormalTok{, }\StringTok{"\#34A853"}\NormalTok{),  }\CommentTok{\# Google colors}
  \FunctionTok{c}\NormalTok{(}\StringTok{"\#1F77B4"}\NormalTok{, }\StringTok{"\#FF7F0E"}\NormalTok{, }\StringTok{"\#2CA02C"}\NormalTok{, }\StringTok{"\#D62728"}\NormalTok{),  }\CommentTok{\# Tableau colors}
  \FunctionTok{c}\NormalTok{(}\StringTok{"\#003f5c"}\NormalTok{, }\StringTok{"\#58508d"}\NormalTok{, }\StringTok{"\#bc5090"}\NormalTok{, }\StringTok{"\#ff6361"}\NormalTok{),  }\CommentTok{\# Viridis{-}like}
  \FunctionTok{c}\NormalTok{(}\StringTok{"\#0073C2"}\NormalTok{, }\StringTok{"\#EFC000"}\NormalTok{, }\StringTok{"\#868686"}\NormalTok{, }\StringTok{"\#CD534C"}\NormalTok{),  }\CommentTok{\# IBM colors}
  \FunctionTok{c}\NormalTok{(}\StringTok{"\#7F3C8D"}\NormalTok{, }\StringTok{"\#11A579"}\NormalTok{, }\StringTok{"\#3969AC"}\NormalTok{, }\StringTok{"\#F2B701"}\NormalTok{)   }\CommentTok{\# Colorbrewer}
\NormalTok{)}
\NormalTok{paleta\_seleccionada }\OtherTok{\textless{}{-}} \FunctionTok{sample}\NormalTok{(paletas\_colores, }\DecValTok{1}\NormalTok{)[[}\DecValTok{1}\NormalTok{]]}
\NormalTok{color\_encabezado }\OtherTok{\textless{}{-}}\NormalTok{ paleta\_seleccionada[}\DecValTok{1}\NormalTok{]}
\NormalTok{color\_primer\_puesto }\OtherTok{\textless{}{-}}\NormalTok{ paleta\_seleccionada[}\DecValTok{2}\NormalTok{]}
\NormalTok{color\_segundo\_puesto }\OtherTok{\textless{}{-}}\NormalTok{ paleta\_seleccionada[}\DecValTok{3}\NormalTok{]}
\NormalTok{color\_tercer\_puesto }\OtherTok{\textless{}{-}}\NormalTok{ paleta\_seleccionada[}\DecValTok{4}\NormalTok{]}
\end{Highlighting}
\end{Shaded}

\begin{Shaded}
\begin{Highlighting}[]
\CommentTok{\# Crear código Python para generar la tabla usando matplotlib}
\NormalTok{codigo\_python }\OtherTok{\textless{}{-}} \FunctionTok{paste0}\NormalTok{(}\StringTok{\textquotesingle{}}
\StringTok{import matplotlib}
\StringTok{matplotlib.use("Agg")}
\StringTok{import matplotlib.pyplot as plt}
\StringTok{import matplotlib.patches as patches}
\StringTok{from matplotlib.gridspec import GridSpec}
\StringTok{import numpy as np}

\StringTok{\# Configurar colores}
\StringTok{color\_encabezado = "\textquotesingle{}}\NormalTok{, color\_encabezado, }\StringTok{\textquotesingle{}"}
\StringTok{color\_primer\_puesto = "\textquotesingle{}}\NormalTok{, color\_primer\_puesto, }\StringTok{\textquotesingle{}"}
\StringTok{color\_segundo\_puesto = "\textquotesingle{}}\NormalTok{, color\_segundo\_puesto, }\StringTok{\textquotesingle{}"}
\StringTok{color\_tercer\_puesto = "\textquotesingle{}}\NormalTok{, color\_tercer\_puesto, }\StringTok{\textquotesingle{}"}

\StringTok{\# Crear figura y configurar grid}
\StringTok{fig = plt.figure(figsize=(8, 3.5))}
\StringTok{gs = GridSpec(4, 2, height\_ratios=[1, 1, 1, 1], width\_ratios=[3, 7])}

\StringTok{\# Encabezado (ocupa todo el ancho)}
\StringTok{ax\_header = plt.subplot(gs[0, :])}
\StringTok{ax\_header.set\_facecolor(color\_encabezado)}
\StringTok{ax\_header.text(0.5, 0.5, "Distribución del premio de \textquotesingle{}}\NormalTok{, premio\_total, }\StringTok{\textquotesingle{} millones de pesos",}
\StringTok{              ha="center", va="center", color="white", fontweight="bold", fontsize=11)}
\StringTok{ax\_header.set\_xticks([])}
\StringTok{ax\_header.set\_yticks([])}
\StringTok{ax\_header.spines["top"].set\_visible(True)}
\StringTok{ax\_header.spines["bottom"].set\_visible(True)}
\StringTok{ax\_header.spines["left"].set\_visible(True)}
\StringTok{ax\_header.spines["right"].set\_visible(True)}

\StringTok{\# Primera fila}
\StringTok{ax\_p1\_label = plt.subplot(gs[1, 0])}
\StringTok{ax\_p1\_label.set\_facecolor(color\_primer\_puesto)}
\StringTok{ax\_p1\_label.text(0.5, 0.5, "Primer puesto", ha="center", va="center", color="white", fontweight="bold", fontsize=10)}
\StringTok{ax\_p1\_label.set\_xticks([])}
\StringTok{ax\_p1\_label.set\_yticks([])}
\StringTok{ax\_p1\_label.spines["top"].set\_visible(True)}
\StringTok{ax\_p1\_label.spines["bottom"].set\_visible(True)}
\StringTok{ax\_p1\_label.spines["left"].set\_visible(True)}
\StringTok{ax\_p1\_label.spines["right"].set\_visible(True)}

\StringTok{ax\_p1\_value = plt.subplot(gs[1, 1])}
\StringTok{ax\_p1\_value.set\_facecolor("white")}
\StringTok{ax\_p1\_value.text(0.5, 0.5, "\textquotesingle{}}\NormalTok{, fraccion\_primer\_puesto, }\StringTok{\textquotesingle{} del \textquotesingle{}}\NormalTok{, termino\_dinero, }\StringTok{\textquotesingle{} total",}
\StringTok{                ha="center", va="center", color="black", fontsize=10)}
\StringTok{ax\_p1\_value.set\_xticks([])}
\StringTok{ax\_p1\_value.set\_yticks([])}
\StringTok{ax\_p1\_value.spines["top"].set\_visible(True)}
\StringTok{ax\_p1\_value.spines["bottom"].set\_visible(True)}
\StringTok{ax\_p1\_value.spines["left"].set\_visible(True)}
\StringTok{ax\_p1\_value.spines["right"].set\_visible(True)}

\StringTok{\# Segunda fila}
\StringTok{ax\_p2\_label = plt.subplot(gs[2, 0])}
\StringTok{ax\_p2\_label.set\_facecolor(color\_segundo\_puesto)}
\StringTok{ax\_p2\_label.text(0.5, 0.5, "Segundo puesto", ha="center", va="center", color="white", fontweight="bold", fontsize=10)}
\StringTok{ax\_p2\_label.set\_xticks([])}
\StringTok{ax\_p2\_label.set\_yticks([])}
\StringTok{ax\_p2\_label.spines["top"].set\_visible(True)}
\StringTok{ax\_p2\_label.spines["bottom"].set\_visible(True)}
\StringTok{ax\_p2\_label.spines["left"].set\_visible(True)}
\StringTok{ax\_p2\_label.spines["right"].set\_visible(True)}

\StringTok{ax\_p2\_value = plt.subplot(gs[2, 1])}
\StringTok{ax\_p2\_value.set\_facecolor("white")}
\StringTok{ax\_p2\_value.text(0.5, 0.5, "\textquotesingle{}}\NormalTok{, fraccion\_segundo\_puesto, }\StringTok{\textquotesingle{} del \textquotesingle{}}\NormalTok{, termino\_dinero, }\StringTok{\textquotesingle{} total",}
\StringTok{                ha="center", va="center", color="black", fontsize=10)}
\StringTok{ax\_p2\_value.set\_xticks([])}
\StringTok{ax\_p2\_value.set\_yticks([])}
\StringTok{ax\_p2\_value.spines["top"].set\_visible(True)}
\StringTok{ax\_p2\_value.spines["bottom"].set\_visible(True)}
\StringTok{ax\_p2\_value.spines["left"].set\_visible(True)}
\StringTok{ax\_p2\_value.spines["right"].set\_visible(True)}

\StringTok{\# Tercera fila}
\StringTok{ax\_p3\_label = plt.subplot(gs[3, 0])}
\StringTok{ax\_p3\_label.set\_facecolor(color\_tercer\_puesto)}
\StringTok{ax\_p3\_label.text(0.5, 0.5, "Tercer puesto", ha="center", va="center", color="white", fontweight="bold", fontsize=10)}
\StringTok{ax\_p3\_label.set\_xticks([])}
\StringTok{ax\_p3\_label.set\_yticks([])}
\StringTok{ax\_p3\_label.spines["top"].set\_visible(True)}
\StringTok{ax\_p3\_label.spines["bottom"].set\_visible(True)}
\StringTok{ax\_p3\_label.spines["left"].set\_visible(True)}
\StringTok{ax\_p3\_label.spines["right"].set\_visible(True)}

\StringTok{ax\_p3\_value = plt.subplot(gs[3, 1])}
\StringTok{ax\_p3\_value.set\_facecolor("white")}
\StringTok{ax\_p3\_value.text(0.5, 0.5, "el \textquotesingle{}}\NormalTok{, termino\_dinero, }\StringTok{\textquotesingle{} restante",}
\StringTok{                ha="center", va="center", color="black", fontsize=10)}
\StringTok{ax\_p3\_value.set\_xticks([])}
\StringTok{ax\_p3\_value.set\_yticks([])}
\StringTok{ax\_p3\_value.spines["top"].set\_visible(True)}
\StringTok{ax\_p3\_value.spines["bottom"].set\_visible(True)}
\StringTok{ax\_p3\_value.spines["left"].set\_visible(True)}
\StringTok{ax\_p3\_value.spines["right"].set\_visible(True)}

\StringTok{plt.tight\_layout(pad=0.5)}

\StringTok{\# Guardar la figura}
\StringTok{plt.savefig("tabla\_distribucion.png", dpi=150, bbox\_inches="tight")}
\StringTok{plt.savefig("tabla\_distribucion.pdf", dpi=150, bbox\_inches="tight")}
\StringTok{plt.close()}
\StringTok{\textquotesingle{}}\NormalTok{)}

\CommentTok{\# Ejecutar código Python para generar la figura}
\FunctionTok{py\_run\_string}\NormalTok{(codigo\_python)}
\end{Highlighting}
\end{Shaded}

\section{Question}\label{question}

En un(a) comunidad se realizará un(a) olimpiada deportiva para menores
de 15 años. Para incentivar a los(las) deportistas, la comunidad cuenta
con 40 millones de pesos, que se repartirán entre los(as) tres primeros
lugares, como se indica a continuación:

\includegraphics[width=0.9\linewidth,height=\textheight,keepaspectratio]{tabla_distribucion.png}

¿Qué cantidad de premio recibe el tercer puesto?

\subsection{Answerlist}\label{answerlist}

\begin{itemize}
\tightlist
\item
  2 millones.
\item
  12 millones.
\item
  5 millones.
\item
  8 millones.
\end{itemize}

\section{Solution}\label{solution}

Para resolver este problema, debemos calcular qué fracción del premio
total corresponde al tercer puesto y luego convertirlo a millones de
pesos.

\subsubsection{Paso 1: Identificar los datos del
problema}\label{paso-1-identificar-los-datos-del-problema}

\begin{itemize}
\tightlist
\item
  Premio total: 40 millones de pesos
\item
  Primer puesto: 2/3 del dinero total
\item
  Segundo puesto: 1/5 del dinero total
\item
  Tercer puesto: el dinero restante
\end{itemize}

\subsubsection{Paso 2: Convertir las fracciones a
decimales}\label{paso-2-convertir-las-fracciones-a-decimales}

\begin{itemize}
\tightlist
\item
  Primer puesto: 2/3 = 0.6666667
\item
  Segundo puesto: 1/5 = 0.2
\end{itemize}

\subsubsection{Paso 3: Calcular la fracción que corresponde al tercer
puesto}\label{paso-3-calcular-la-fracciuxf3n-que-corresponde-al-tercer-puesto}

Para calcular la fracción del tercer puesto, restamos del total (1) las
fracciones del primer y segundo puesto:

\begin{itemize}
\tightlist
\item
  Fracción del tercer puesto = 1 - (0.6666667 + 0.2)
\item
  Fracción del tercer puesto = 1 - 0.8666667
\item
  Fracción del tercer puesto = 0.1333333
\end{itemize}

\subsubsection{Paso 4: Calcular el monto en millones de pesos para el
tercer
puesto}\label{paso-4-calcular-el-monto-en-millones-de-pesos-para-el-tercer-puesto}

Multiplicamos la fracción del tercer puesto por el premio total:

\begin{itemize}
\tightlist
\item
  Monto del tercer puesto = 0.1333333 × 40 millones
\item
  Monto del tercer puesto = 5 millones de pesos
\end{itemize}

\subsubsection{Verificación}\label{verificaciuxf3n}

Comprobemos que la suma de los tres montos es igual al premio total:

\begin{itemize}
\tightlist
\item
  Primer puesto: 27 millones
\item
  Segundo puesto: 8 millones
\item
  Tercer puesto: 5 millones
\item
  Total: 40 millones
\end{itemize}

Como 40 = 40, confirmamos que nuestra respuesta es correcta.

Por lo tanto, el tercer puesto recibe 5 millones de pesos.

\subsection{Answerlist}\label{answerlist-1}

\begin{itemize}
\tightlist
\item
  Falso
\item
  Falso
\item
  Verdadero
\item
  Falso
\end{itemize}

\section{Meta-information}\label{meta-information}

exname: fracciones\_reparto\_premio extype: schoice exsolution: 0010
exshuffle: TRUE exsection: Aritmética\textbar Fracciones\textbar Reparto
proporcional

\end{document}
